\documentclass{article}
\usepackage{amsmath}
\usepackage{graphicx}

\title{RuddyDoc LaTeX Test Document}
\author{RuddyDoc Test Suite}
\date{April 2024}

\begin{document}

\maketitle

\section{Introduction}

This is the introduction section with regular text. It demonstrates basic paragraph parsing in \LaTeX{} documents.

This is a second paragraph with \textbf{bold text}, \emph{emphasized text}, and \texttt{monospace text}.

\section{Lists and Structure}

\subsection{Unordered Lists}

Here is an unordered list:

\begin{itemize}
    \item First item in the list
    \item Second item with some \textbf{formatting}
    \item Third item
\end{itemize}

\subsection{Ordered Lists}

And here is an ordered list:

\begin{enumerate}
    \item Step one
    \item Step two
    \item Step three
\end{enumerate}

\subsection{Nested Lists}

\begin{itemize}
    \item Outer item one
    \begin{itemize}
        \item Inner item one
        \item Inner item two
    \end{itemize}
    \item Outer item two
\end{itemize}

\section{Tables}

\begin{table}[h]
\centering
\begin{tabular}{|l|c|r|}
\hline
\textbf{Column A} & \textbf{Column B} & \textbf{Column C} \\
\hline
Left aligned & Centered & Right aligned \\
Row 2, Col 1 & Row 2, Col 2 & Row 2, Col 3 \\
Row 3, Col 1 & Row 3, Col 2 & Row 3, Col 3 \\
\hline
\end{tabular}
\caption{A sample table with three columns}
\label{tab:sample}
\end{table}

\section{Mathematics}

Inline math: The quadratic formula is $x = \frac{-b \pm \sqrt{b^2-4ac}}{2a}$.

Display math:
\begin{equation}
E = mc^2
\label{eq:einstein}
\end{equation}

As shown in equation~\ref{eq:einstein}, energy and mass are related.

\section{Figures}

\begin{figure}[h]
\centering
\includegraphics[width=0.5\textwidth]{images/plot.png}
\caption{A sample plot showing data trends}
\label{fig:plot}
\end{figure}

Figure~\ref{fig:plot} shows the data visualization.

\section{Code Blocks}

Here is a verbatim code block:

\begin{verbatim}
def factorial(n):
    if n <= 1:
        return 1
    return n * factorial(n-1)
\end{verbatim}

\section{Footnotes and References}

This text has a footnote\footnote{This is the footnote text explaining something.} that provides additional context.

We can reference Table~\ref{tab:sample} and Figure~\ref{fig:plot} using labels.

% This is a comment and should not appear in the parsed output

\subsubsection{A Third-Level Heading}

This demonstrates a third-level heading structure.

\section{Conclusion}

This concludes the LaTeX test document with all major structural elements covered.

\end{document}
